\chapter{Nullelemente aus Potenzhalbgruppen erkennen}
\hypertarget{zero-transport.reading}{}

Ein Isomorphismus zwischen den nichtleeren Potenzhalbgruppen zweier
Halbgruppen erhält die Existenz eines Nullelements. Das gilt auch für
unendliche Halbgruppen und benötigt keine Einschränkung der Zahl ihrer
Produktwerte. Genauer gilt
\FormulaRefAuto{PowerSemigroupZeroTransport}[thm]:
Sind \(S,T\) Halbgruppen, ist \(0_S\) eine Null von \(S\) und ist
\[
 \Phi:\PowNE S\longrightarrow\PowNE T
\]
ein Halbgruppenisomorphismus, dann besitzt \(T\) eine Null \(0_T\), und
\[
 \Phi(\{0_S\})=\{0_T\}.
\]
Eine Null ist ein Element, das bei Multiplikation von beiden Seiten
absorbiert: \(0_Ss=s0_S=0_S\) für jedes \(s\in S\).
Die Definition und die Eindeutigkeit dieses Elements stehen in
\FormulaRefAuto{SemigroupWithZeroDef}[def] und
\FormulaRefAuto{SemigroupZeroUnique}[thm] aus Band~33.
Die Ergänzung verwendet außerdem die Potenzhalbgruppen und
Isomorphieregeln aus Band~28. Ihre Beweistabellen tragen die eigene
Nummerierung~28Z; ältere Satznummern bleiben dadurch unverändert.

\section{Warum eine Null der Potenzhalbgruppe noch nicht genügt}

Die Elemente von \(\PowNE S\) sind die nichtleeren Teilmengen von \(S\).
Ihr Produkt ist
\[
 AB=\{ab\mid a\in A,\ b\in B\}.
\]
Weil \(0_S\) jedes Grundelement absorbiert, absorbiert die Einermenge
\(\{0_S\}\) jede nichtleere Teilmenge:
\[
 \{0_S\}A=A\{0_S\}=\{0_S\}\qquad(A\in\PowNE S).
\]
Setzen wir
\[
 Z=\Phi(\{0_S\}),
\]
so ist \(Z\) eine nichtleere Teilmenge von \(T\) und die Null
der Potenzhalbgruppe \(\PowNE T\). Tatsächlich besitzt jedes
\(Y\in\PowNE T\) wegen der Surjektivität von \(\Phi\) ein Urbild
\(A\in\PowNE S\). Daher ist
\[
 ZY=\Phi(\{0_S\})\Phi(A)
    =\Phi(\{0_S\}A)=\Phi(\{0_S\})=Z,
\]
und entsprechend \(YZ=Z\).

Damit ist bisher eine \emph{Menge} als absorbierendes Element der
Potenzhalbgruppe gefunden. Noch ist nicht bewiesen, dass diese Menge
nur ein Grundelement enthält. Genau das ist der verbleibende Schritt:
Wir müssen \(Z=\{0_T\}\) für ein \(0_T\in T\) zeigen. Erst dann
folgt aus \(Z\{t\}=\{t\}Z=Z\), dass
\(0_Tt=t0_T=0_T\).

\section{Eine Menge aus Urbildern von Einermengen}

Für jedes \(z\in Z\) ist \(\{z\}\) ein Element von \(\PowNE T\).
Die Umkehrabbildung \(\Phi^{-1}\) ist auf dieser Einermenge definiert.
Wir schreiben
\[
 A_z=\Phi^{-1}(\{z\}),\qquad
 W=\bigcup_{z\in Z}A_z.
\]
Dabei ist \(A_z\) ein einzelner Wert der Umkehrabbildung, also eine
nichtleere Teilmenge von \(S\); es handelt sich nicht um das Urbild
einer Mengenfamilie. Die Vereinigung \(W\) sammelt die Grundelemente
aller dieser Teilmengen. Jedes ihrer Elemente liegt in \(S\), also
\(W\subseteq S\). Weil \(Z\) nicht leer ist, gibt es ein \(z\in Z\);
wegen \(A_z\neq\varnothing\) liegt dann auch ein Element in \(W\).
Somit ist \(W\in\PowNE S\), und \(\Phi(W)\) ist definiert.
Hier werden keine beliebigen Urbilder ausgewählt: Die Bijektion
\(\Phi\) bestimmt jedes \(A_z\) eindeutig.

Die Umkehrabbildung erhält Produkte
(\FormulaRefAuto{SemigroupIsoInverseOperation}[thm]). Deshalb gilt
für \(z,w\in Z\)
\[
 A_zA_w
 =\Phi^{-1}(\{z\})\Phi^{-1}(\{w\})
 =\Phi^{-1}(\{zw\})=A_{zw}.
\]
Der Index \(zw\) liegt wieder in \(Z\), denn
\(\{z\}Z=Z\).

\section{Jedes \texorpdfstring{\(A_z\)}{Az} lässt die Vereinigung fest}

Fixieren wir \(z\in Z\). Aus der Nullwirkung von \(Z\) folgt
\[
 \{z\}Z=Z.
\]
Elementweise bedeutet dies: Die Produkte \(zw\) mit \(w\in Z\)
durchlaufen ganz \(Z\). Es kann Wiederholungen geben; für eine
Vereinigung spielen Wiederholungen keine Rolle. Folglich ist
\[
 \bigcup_{w\in Z}A_{zw}=\bigcup_{v\in Z}A_v=W.
\]
Dabei werden beide Inklusionen benutzt: Jedes \(zw\) liegt in \(Z\),
und zu jedem \(v\in Z\) gibt es ein \(w\in Z\) mit \(zw=v\).

Das Mengenprodukt verteilt sich über Vereinigungen. Ein Element von
\(A_zW\) hat die Form \(ab\) mit \(a\in A_z\) und \(b\in W\).
Die zweite Zugehörigkeit bedeutet, dass \(b\) in mindestens einem
\(A_w\) liegt. Somit liegt \(ab\) in einem \(A_zA_w\).
Umgekehrt ist jedes Element eines solchen Produkts auch in \(A_zW\).
Also erhalten wir
\[
 A_zW
 =A_z\Bigl(\bigcup_{w\in Z}A_w\Bigr)
 =\bigcup_{w\in Z}A_zA_w
 =\bigcup_{w\in Z}A_{zw}=W.
\]
Die allgemeinen Vereinigungsrechnungen werden in
\FormulaRefAuto{ZeroTransportProductUnion}[thm] und
\FormulaRefAuto{ZeroTransportReindexedUnion}[thm] vollständig bewiesen.

\section{Das Bild der Vereinigung ist die Potenznull}

Nun wenden wir \(\Phi\) auf die Gleichheit \(A_zW=W\) an.
Da \(\Phi(A_z)=\{z\}\), ergibt die Produktverträglichkeit
\[
 \{z\}\Phi(W)
 =\Phi(A_z)\Phi(W)
 =\Phi(A_zW)=\Phi(W)
 \qquad(z\in Z).
\]
Jedes Grundelement \(z\in Z\) lässt also die Menge \(\Phi(W)\)
bei Linksmultiplikation unverändert. Nehmen wir die Vereinigung dieser
gleichen Mengen, erhalten wir
\[
 Z\Phi(W)
 =\bigcup_{z\in Z}\{z\}\Phi(W)
 =\bigcup_{z\in Z}\Phi(W)=\Phi(W).
\]
Im letzten Schritt wird erneut \(Z\neq\varnothing\) gebraucht:
Die Vereinigung einer nichtleeren Familie, deren sämtliche Werte
dieselbe Menge sind, ist diese Menge.
Andererseits ist \(Z\) die Null der Potenzhalbgruppe. Da
\(\Phi(W)\in\PowNE T\), ist daher auch
\[
 Z\Phi(W)=Z.
\]
Die beiden Rechnungen haben dieselbe linke Seite und liefern
\(\Phi(W)=Z=\Phi(\{0_S\})\). Aus der Injektivität folgt
\[
 W=\{0_S\}.
\]

In diesem Argument wird \(\Phi\) nur auf Produkte und Gleichheiten
angewandt. Die beiden Vereinigungen werden jeweils innerhalb einer
Potenzhalbgruppe berechnet. Eine Erhaltung von Vereinigungen durch
\(\Phi\) ist nicht vorausgesetzt und wurde auch nicht benutzt.

\section{Von der Einermenge zum Nullelement}

Wählen wir ein \(z\in Z\). Die Menge \(A_z\) ist nicht leer und
liegt in ihrer Vereinigung \(W=\{0_S\}\). Die einzige nichtleere
Teilmenge dieser Einermenge ist die Einermenge selbst
(\FormulaRefAuto{NonemptySubsetSingleton}[thm]). Daher ist
\[
 A_z=\{0_S\},\qquad
 \{z\}=\Phi(A_z)=\Phi(\{0_S\})=Z.
\]
Damit besitzt \(Z\) tatsächlich genau ein Element. Setzen wir
\(0_T=z\), so liefert für jedes \(t\in T\) die Nullwirkung von
\(Z\)
\[
 \{0_Tt\}=\{0_T\}\{t\}=Z\{t\}=Z=\{0_T\},
\]
und ebenso \(\{t0_T\}=\{0_T\}\). Die Gleichheit zweier
Einermengen erzwingt die Gleichheit ihrer Elemente. Also ist
\(0_Tt=t0_T=0_T\), wie behauptet.

Der Beweis verwendet keine Endlichkeit und kein Zählargument.
Er zeigt genau die Erhaltung der Null-Einermenge. Die Bilder anderer
Einermengen werden dadurch nicht festgelegt.
Die vollständige formale Schlussableitung steht in
\FormulaRefAuto{PowerSemigroupZeroTransport}[thm] der
\href{_B28-zero-transport-proofs.pdf\#zero-transport.proofs}{zugehörigen Beweistabellen}.
